\section{Evaluation}
\label{sec:evaluation}

\begin{itemize}
\item[\textbf{RQ1}] Under crashes, does AEP eliminate \emph{undetected}
      duplicates, and what does its residual ambiguity depend on?
\item[\textbf{RQ2}] Which of the protocol's two durability mechanisms produces
      which guarantee --- and against which faults is each one observable?
\item[\textbf{RQ3}] What does the protocol cost, and how much of that cost is
      a choice?
\item[\textbf{RQ4}] How does recovery behave?
\end{itemize}

\subsection{Setup}
\label{sec:eval-setup}

\paragraph{Systems} The seven of \cref{tab:trilemma}, implemented as thin
variants over one connector, one workload driver and one oracle, so that only
the differences differ. Each carries a machine-readable descriptor that tests
check against the implementation, and a crash point that does not exist in a
given system is marked inapplicable and never run rather than silently skipped.

\paragraph{Faults} Workers are separate OS processes killed with real
\texttt{SIGKILL} at one of the six crash points of \cref{tab:crashpoints}. A
\emph{regime} is a named fault condition --- not a matrix dimension --- and we
report each separately because they are different experiments:

\begin{itemize}
\item \textbf{crashed}: every execution is killed at the cell's crash point.
      This is the regime RQ1 is about.
\item \textbf{crash-free} (\texttt{p0}): nothing is injected. The only regime
      RQ3 may use.
\item \textbf{\texttt{redis-kill-preack}}: no worker crash; Redis is hard-killed
      with \texttt{docker kill -s KILL} at the instruction boundary between the
      intent CAS and the barrier acknowledgement, then restarted and verified.
      This is RQ2's ablation.
\end{itemize}

\paragraph{Oracle} A standalone mock provider records every \emph{actually
applied} mutation in a SQLite ledger (WAL, \texttt{synchronous=FULL}, one
transaction per mutation). Duplicates are counted from that ledger, never from
the system's own opinion of itself. The analysis is permitted to read exactly
two files per run --- the event log and the oracle ledger --- and a source gate
that parses the analysis modules fails the build if they import anything that
could reach live protocol state.

\paragraph{Statistics} Rates are over executions. Per-cell rate intervals use
a seeded run-cluster bootstrap (\BootstrapResamples{} resamples); the
B3--AEP ambiguity difference uses a stratified run-cluster bootstrap that
preserves crash-point, endpoint-capability, and keying strata. Fisher's exact
tests compare execution counts and therefore do not adjust for run clustering;
we report them descriptively, not as evidence of equivalence. Seeds and
repetition counts are recorded per run.

\paragraph{Two reporting rules we hold ourselves to} First, \emph{no pooled
table}. A rate computed across fault regimes is a property of how many runs of
each kind were collected, not of any system; our own analysis tool prints a
warning whenever more than one regime is present, and every number below names
its regime. Second, \emph{no absolute timing from an ungated host}. A run
contributes a duration only if the operator declared before the run that the
host cannot suspend \emph{and} the run's wall-versus-monotonic divergence stays
within tolerance. Both halves are necessary: a host that merely was not idle
long enough to suspend is indistinguishable from one that cannot. Under this
gate, \emph{zero} runs collected before this rule existed contribute a timing
number, and none does. Both rules were adopted after the first collection ---
the pooling rule when a mid-study review found the pooled summary table
unusable as a source, the timing gate after a host suspension was discovered
mid-run --- and every rate and duration below is reported under them.

\subsection{RQ1 --- undetected duplicates and the shape of the residual}
\label{sec:eval-rq1}

\Cref{tab:outcomes} is the anchor result; the supplementary material shows the
same executions pooled to one bar pair per system.

% GENERATED by scripts/paper_tables.py -- do not edit by hand.
% Source: analysis/per-cell-metrics.csv
% Filter: regime == '(session-3)' ONLY.
%   Every execution in these cells was killed at the cell's crash
%   point. The crash-free (p0) and hard-Redis-kill (redis-kill-preack)
%   regimes are DIFFERENT EXPERIMENTS and are reported separately;
%   pooling them is what got analysis/table-1.csv banned.
% Each cell = sum(successes)/sum(total) over the six crash points.
\begin{table*}[t]
\centering
\caption{The trilemma, measured. Every execution in every cell was killed at one of the six crash points of \cref{tab:crashpoints}; rates are over executions, pooled across crash points within one endpoint capability. \textsc{auth}/\textsc{pos-only}/\textsc{none} are the reconciliation capabilities of \cref{tab:capabilities}. AEP-full is the only system with a nonzero declared-ambiguity column, and the only one whose other two columns are zero everywhere. Source: \texttt{per-cell-metrics.csv}, crashed regime only.}
\label{tab:outcomes}
\small
\begin{tabular}{@{}lccccccccc@{}}
\toprule
& \multicolumn{3}{c}{undetected duplicate}& \multicolumn{3}{c}{lost effect}& \multicolumn{3}{c}{declared ambiguity}\\
\cmidrule(lr){2-4}\cmidrule(lr){5-7}\cmidrule(lr){8-10}
system & \textsc{auth} & \textsc{pos-only} & \textsc{none} & \textsc{auth} & \textsc{pos-only} & \textsc{none} & \textsc{auth} & \textsc{pos-only} & \textsc{none}\\
\midrule
B0 naive retry & 0.8200 & 0.8000 & 0.7733 & 0.0133 & 0.0133 & 0.0067 & 0.0000 & 0.0000 & 0.0000\\
B1 lease-only & 0.8133 & 0.8067 & 0.8267 & 0.0000 & 0.0067 & 0.0067 & 0.0000 & 0.0000 & 0.0000\\
B2 CAS-only & 0.8067 & 0.7933 & 0.7800 & 0.0067 & 0.0067 & 0.0000 & 0.0000 & 0.0000 & 0.0000\\
B4 durable, $\infty$ attempts & 0.9500 & 0.5889 & 0.5611 & 0.0000 & 0.0000 & 0.0111 & 0.0000 & 0.0000 & 0.0000\\
B4b durable, 1 attempt & --- & 0.0000 & 0.0000 & --- & 0.5056 & 0.5167 & --- & 0.0000 & 0.0000\\
B3 intent, no barrier & 0.0000 & 0.0000 & 0.0000 & 0.0000 & 0.0000 & 0.0000 & 0.0000 & 0.3667 & 0.7167\\
\textbf{AEP-full} & 0.0000 & 0.0000 & 0.0000 & 0.0000 & 0.0000 & 0.0000 & 0.0000 & 0.3500 & 0.7222\\
\bottomrule
\end{tabular}
\end{table*}

% Denominators (executions) per (system, response class):
%   B0_NAIVE_RETRY                     AUTHORITATIVE_READBACK   executions=150   runs=15   crash_points=5
%   B0_NAIVE_RETRY                     POSITIVE_ONLY_READBACK   executions=150   runs=15   crash_points=5
%   B0_NAIVE_RETRY                     NO_READBACK              executions=150   runs=15   crash_points=5
%   B1_LEASE_ONLY                      AUTHORITATIVE_READBACK   executions=150   runs=15   crash_points=5
%   B1_LEASE_ONLY                      POSITIVE_ONLY_READBACK   executions=150   runs=15   crash_points=5
%   B1_LEASE_ONLY                      NO_READBACK              executions=150   runs=15   crash_points=5
%   B2_CAS_ONLY                        AUTHORITATIVE_READBACK   executions=150   runs=15   crash_points=5
%   B2_CAS_ONLY                        POSITIVE_ONLY_READBACK   executions=150   runs=15   crash_points=5
%   B2_CAS_ONLY                        NO_READBACK              executions=150   runs=15   crash_points=5
%   B4_DURABLE_WORKFLOW                AUTHORITATIVE_READBACK   executions=20    runs=2    crash_points=1
%   B4_DURABLE_WORKFLOW                POSITIVE_ONLY_READBACK   executions=180   runs=18   crash_points=6
%   B4_DURABLE_WORKFLOW                NO_READBACK              executions=180   runs=18   crash_points=6
%   B4B_DURABLE_WORKFLOW_AT_MOST_ONCE  POSITIVE_ONLY_READBACK   executions=180   runs=18   crash_points=6
%   B4B_DURABLE_WORKFLOW_AT_MOST_ONCE  NO_READBACK              executions=180   runs=18   crash_points=6
%   B3_INTENT_NO_BARRIER               AUTHORITATIVE_READBACK   executions=180   runs=18   crash_points=6
%   B3_INTENT_NO_BARRIER               POSITIVE_ONLY_READBACK   executions=180   runs=18   crash_points=6
%   B3_INTENT_NO_BARRIER               NO_READBACK              executions=180   runs=18   crash_points=6
%   AEP_FULL                           AUTHORITATIVE_READBACK   executions=180   runs=18   crash_points=6
%   AEP_FULL                           POSITIVE_ONLY_READBACK   executions=180   runs=18   crash_points=6
%   AEP_FULL                           NO_READBACK              executions=180   runs=18   crash_points=6


Three readings, in the order they matter.

\paragraph{AEP's two silent columns are empty, everywhere measured} Across
\AepExecAuth{}, \AepExecPosOnly{} and \AepExecNoReadback{} crashed executions on
the three capability classes, AEP-full recorded an undetected-duplicate rate of
\AepDupAuth{}, \AepDupPosOnly{} and \AepDupNoReadback{}, and a lost-effect rate
of \AepLostAuth{}, \AepLostPosOnly{} and \AepLostNoReadback{}. Every one of
those executions was killed by \texttt{SIGKILL} at a named point around the
external call. The matrix runner halted on this condition: a single undetected
duplicate in AEP-full stops the matrix, a rule carried in the runner before any
of these cells were collected. We do not call it a pre-registration. The
governing document at the time described a contrary result as a defect to fix
rather than a finding to report, and this project's pre-registration practice
began later; what the runner gives is a stopping rule fixed in advance, not a
commitment about how the outcome would have been written up. It did not occur.

\paragraph{The baselines without a pre-dispatch record duplicate in most
crashed executions} B0, B1 and B2 sit between \BaselineDupLow{} and
\BaselineDupHigh{}, and the weakest of the three comparisons against AEP-full
is significant at $p = \BaselineDupMaxP{}$ by Fisher's exact test.
% \BaselineDupLow / \BaselineDupHigh: generated/numbers.tex, from
% per-cell-metrics.csv over {B0,B1,B2} x all collected response classes.
% p-value source: analysis/comparisons-vs-aep-full.csv, crashed regime.
The gap between B1/B2 and B0 is small and that is the finding: a lease and a
fenced CAS are doing exactly what they promise and neither promise is about the
external call (\cref{fig:bycrashpoint}).

\paragraph{A durable log is not enough, and the two configurations of one
engine land on two different corners} B4 has a write-ahead record,
acknowledged durable by the same \texttt{WAITAOF} barrier AEP uses. It is not
ablated on durability. It duplicates at \BfourDupPosOnly{} on
\textsc{pos-only} and \BfourDupNoReadback{} on \textsc{none}
(\BfourExecPosOnly{} and \BfourExecNoReadback{} executions), and it duplicates
on \textsc{auth} as well, at \BfourDupAuth{} over \BfourExecAuth{}
executions. The endpoint's
reconciliation capability does not rescue it, which is the point --- a
read-back tells a recovering workflow what the provider currently holds, and
B4's duplicate is a second application it chose to make, not a fact it failed
to look up.

B4b is the same engine with Maximum Attempts set to $1$, and it is the cleaner
half of the demonstration: over \BfourbExecPosOnly{} and
\BfourbExecNoReadback{} executions it records \emph{zero} undetected
duplicates on both classes and loses the effect in \BfourbLostPosOnly{} and
\BfourbLostNoReadback{} of them instead, and \BfourbLostAuth{} over
\BfourbExecAuth{} on \textsc{auth}. One configuration line moves
the engine from one silent corner of \cref{sec:trilemma} to the other, and
neither configuration reaches the third: across all \BfourFamilyAmbCells{}
B4 and B4b cells the declared-ambiguity rate never exceeds
\BfourFamilyAmbMax{}.

\paragraph{The engine itself was run, and it separates what B4 licenses from
what it does not} We ran the production engine against the same non-cooperative
endpoint, in the two configurations B4 and B4b instantiate, at the one crash
point both the engine and our model can be cut at. With unlimited attempts it
duplicated at \BfiveDup{} over \BfiveExec{} executions and lost effects at
\BfiveLost{}; with Maximum Attempts set to one it lost at \BfivebLost{} over
\BfivebExec{} executions and duplicated at \BfivebDup{}. Neither configuration
declared ambiguity in any of \BfiveRuns{} non-void runs (\BfiveAmb{} and
\BfivebAmb{} declarations) --- the third corner of \cref{sec:trilemma} going
unreached by the engine itself and not merely by our model of it.

\emph{The direction transfers; the magnitudes do not.} Each configuration lands
in the corner the supplementary comparison assigns it, with exactly zero events
of the other kind. But at that same crash point B4 duplicates at
\BfourAtBarrier{} over \BfourAtBarrierExec{} executions and B4b loses at
\BfourbAtBarrier{} over \BfourbAtBarrierExec{}, roughly six times the engine's
rates. \textbf{B4's rates are our model's, not the product's}, and no rate in
this paper should be read as an engine's.

Four things bound this and none is removed by having run the engine. The
comparison is \emph{incomplete}: the engine has no position answering to
\textsc{after-intent-before-barrier}, because its pre-dispatch record is written
by the server inside one RPC, so there is no worker-side instant between the
write and its acknowledgement --- a hole exactly where B4 and B4b have data at
all three capability classes. The rates are conditional on a
\BfiveStartToClose{}\,ms Start-To-Close timeout that B4 has no analogue for; a
different timeout is a different cell, and no timing number here is any engine's
recovery latency. The interval arithmetic is asymmetric --- B4's intervals rest
on three runs against the engine's thirty, and three clusters put resampled
means on multiples of one thirtieth --- so the separation should be read from
the point estimates, which differ by roughly $0.8$, and not from the intervals
failing to overlap. And the crash window at the dispatch point is sized by the
server and by loopback rather than by the protocol, so its width is not
like-for-like with B4's.

That is the sharpest form of the argument the paper can make \emph{against our
model of a durable-execution engine} --- and that qualification is load-bearing
rather than modest, and the paragraphs above are why: B4 is our implementation
of the pattern, not a product, and we registered before collecting B5 that a
gap of that size would oblige us to re-scope every B4 claim. The shape of the
argument transfers and the magnitude is our model's. The argument is nonetheless
made against a system that has everything the write-ahead pattern is supposed to
provide. The record is necessary and it is not sufficient, because the
record cannot contain the fact that matters --- whether the provider applied
the effect. What separates B3 and AEP-full from B4 is not durability but the
\emph{absence of an edge back into the dispatching state}: the transition
table has no path that re-enters \textsc{about-to-fire}, checked inside the
atomic script (\cref{fig:statemachine}). What separates them from B4b is that
declining to re-send is not the same as knowing whether to.

\subsubsection{The residual is set by the endpoint, not the protocol}

AEP's declared-ambiguity column is \AepAmbAuth{} on \textsc{auth},
\AepAmbPosOnly{} on \textsc{pos-only} and \AepAmbNoReadback{} on
\textsc{none}. That ordering is not a defect that better engineering would
remove; it is \cref{tab:capabilities} showing through. Where the endpoint can
prove absence, recovery resolves everything and nothing is left ambiguous.
Where it can only confirm presence, absence is uninformative and the honest
answer is ``I do not know''. Where it cannot be queried at all, no amount of
protocol recovers the fact.

\paragraph{The \textsc{pos-only} rate depends on how the read-back is keyed, and
we did not expect it to} The read-back in the cells above is keyed by the
caller's own reference. We pre-registered an alternative keying --- an oracle
fingerprint of the request --- as a sensitivity check, and predicted it would
leave every headline rate inside a $\pm\BthreeVsAepAmbMargin{}$ percentage-point
margin. On \textsc{pos-only} it does not: declared ambiguity is
\KeyingAmbiguityCaller{}\% under the caller's reference and
\KeyingAmbiguityOracle{}\% under the fingerprint, a difference of
\KeyingAmbiguityDelta{} percentage points, and the margin was
\BthreeVsAepAmbMargin{}.
% generated/numbers.tex | ws5-2026-09-10/t2-keying/analysis/
% per-cell-metrics.csv against analysis/per-cell-metrics.csv, metric=
% known_ambiguity_rate, system=AEP_FULL, response_class=POSITIVE_ONLY_READBACK,
% pooled over the six crash points as sum(successes)/sum(total).

The two columns that carry the safety claim are unmoved: undetected duplicates
and lost effects are zero under both keyings. What moves is the residual, and
the consequence is a limit on generality rather than on correctness. **The
\textsc{pos-only} ambiguity rate is a property of the protocol \emph{and} of how
the recovering worker phrases its question**, so the figure quoted above should
be read as this keying's, not as the class's. \textsc{auth} is unaffected,
because a read-back that can prove absence resolves the execution whichever key
it is asked under.

The supplementary material plots the same executions as one bar pair per system,
pooled over crash points and capabilities: the three systems on the right have an
empty left bar, and only AEP-full and its ablation B3 have a right bar at all.
The figure adds no number this section does not state; the bars are Wilson
intervals on pooled counts, and the cluster-aware bootstrap intervals in
\texttt{per-cell-metrics.csv} are wider and are the ones to quote.

The rate decomposes further by crash point, and that is where the protocol's
conservatism becomes visible; the supplementary material gives the full
decomposition across all six crash points and three capability classes.

Two cells carry the argument. At \texttt{before\_intent\_write} the rate is
zero in every capability class: the crash precedes the intent, nothing was
promised, and there is nothing to be ambiguous about. At
\texttt{after\_barrier\_before\_dispatch} it is at its maximum --- and that is
the crash point at which \emph{no effect can possibly exist}, because the
worker died before sending anything. AEP declares ambiguity most often exactly
where it is least warranted by the world. That is the false-positive ambiguity
of \cref{sec:p2}, and its direction is the whole design: the protocol's errors
are toward saying ``I might have done something'' when it did not, never toward
silence when it did. The cells where an effect \emph{does} typically exist ---
\texttt{mid\_dispatch} and \texttt{after\_response\_before\_resolution} under
\textsc{pos-only} --- show low ambiguity, because a positive read-back
confirms the effect and the intent resolves.

\subsubsection{Why the provably-empty cell is not resolved}
\label{sec:eval-provable}

The \texttt{after\_barrier\_before\_dispatch} row invites an obvious question:
if the worker died before sending anything, \emph{the system could know that}.
Why is the outcome ambiguous rather than a confirmed non-event?

Because nothing durable records it. A recovery process reading the store
afterwards sees an authorized, unresolved intent --- exactly what it would see
had the worker died \emph{during} transmission, with the request already on the
wire. The two histories are identical in everything durable, so recovery fails
closed, which is the correct behaviour given what it can observe. \textbf{That
is the negative result, and it is the protocol's, not the implementation's:} the
cell is empty in the world and unresolvable from the store. Closing it is
possible and we did not do it. The supplementary material gives the probe detail
and prices the third barrier that would split the two cases --- a trade that
narrows rather than closes the window, and whose value depends entirely on a
deployment's ratio of crash points, which our uniform matrix deliberately does
not model.

% Single column rather than full width. A float's cost is its area, and for a
% graphic -- unlike a table -- narrowing genuinely shrinks the area rather than
% re-wrapping the same rows taller: measured at 0.29 pages saved, against 0.02
% to 0.04 for the three tables that were tested the same way and left alone.
\begin{figure}[t]
\centering
\includegraphics[width=\columnwidth]{figures/figure-2-duplicates-by-crash-point.pdf}
\caption{Undetected duplicate rate by crash point, crashed regime, pooled over
endpoint capabilities weighted by executions. The baselines' rate is not
concentrated at one crash point: any interruption after the request is on the
wire is enough, which is why a mechanism that acts only \emph{after} the call
cannot help.}
% Generated by experiments/analyze.py write_figures(), regime "(session-3)".
\label{fig:bycrashpoint}
\end{figure}

\subsubsection{The cost of the third corner}

Declared ambiguity is not free, and a paper that presented it as a pure win
would be selling something. Each declared ambiguity is an execution that has
stopped, that will not proceed without a human, and whose resolution requires
someone to establish out-of-band what the endpoint will not say. On a
\textsc{no-readback} endpoint under a crash-heavy regime that is most of the
crashed executions. What we claim is not that this is cheap but that it is
\emph{the honest price of the information the endpoint withholds}, and that
the alternative is not fewer incidents but the same incidents discovered later
by a different department. We have not evaluated whether operators handle
declared ambiguities well, or how much a real deployment would generate; both
are stated as gaps in \cref{sec:threats}.

\subsection{RQ2 --- which mechanism produces which guarantee}
\label{sec:eval-rq2}

AEP contains two durability mechanisms and the literature usually treats them
as one. The first is the \emph{pre-dispatch record}: an intent, written and
fenced, that exists before the external call does. The second is the
\emph{barrier}: the \texttt{WAITAOF} round trip that makes the record's
durability a checked precondition of dispatch authority (\cref{sec:p2}). B3 is
the ablation that separates them --- AEP-full with the barrier removed and
nothing else removed --- and running it against three different fault classes
assigns each guarantee to the mechanism that actually produces it.

The assignment is not the one we expected, and it is the paper's central
structural result.

\subsubsection{Detection is the record's, not the barrier's}
\label{sec:eval-detection}

Everything in \cref{sec:eval-rq1} --- zero undetected duplicates, zero lost
effects, a residual of declared ambiguity shaped by endpoint capability ---
survives the barrier's removal intact.

% GENERATED by scripts/paper_tables.py -- do not edit.
% Source: analysis/per-cell-metrics.csv (rates, regime=crashed)
%         analysis/comparisons-vs-aep-full.csv (regime-labelled p-values).
\begin{table}[t]
\centering
\caption{The barrier ablation on the detection metrics. B3 is AEP-full with the \texttt{WAITAOF} barrier removed and nothing else removed. Crashed regime; rates are over executions, pooled across crash points within one capability class. The p-values are execution-level Fisher tests over all capability classes; they do not account for run clustering and are not used as equivalence evidence.}
\label{tab:ablation}
\small
\begin{tabular}{@{}llccc@{}}
\toprule
metric & capability & AEP-full & B3 & $p$\\
\midrule
\multirow{3}{*}{undetected duplicate} & \textsc{auth} & 0.0000 & 0.0000 & \multirow{3}{*}{$1.00$}\\
 & \textsc{pos-only} & 0.0000 & 0.0000 & \\
 & \textsc{none} & 0.0000 & 0.0000 & \\
\midrule
\multirow{3}{*}{lost effect} & \textsc{auth} & 0.0000 & 0.0000 & \multirow{3}{*}{$1.00$}\\
 & \textsc{pos-only} & 0.0000 & 0.0000 & \\
 & \textsc{none} & 0.0000 & 0.0000 & \\
\midrule
\multirow{3}{*}{declared ambiguity} & \textsc{auth} & 0.0000 & 0.0000 & \multirow{3}{*}{$0.95$}\\
 & \textsc{pos-only} & 0.3500 & 0.3667 & \\
 & \textsc{none} & 0.7222 & 0.7167 & \\
\bottomrule
\end{tabular}
\end{table}


Over \BthreeVsAepN{} crashed executions per arm, B3 and AEP-full show no
observed difference in either silent-failure metric. Both record an
undetected-duplicate count of $\BthreeVsAepDupCount{}$ and a lost-effect count
of $\BthreeVsAepLostCount{}$. Declared ambiguity is
\BthreeAmbCount{}/\BthreeVsAepN{} for B3 and
\AepAmbCount{}/\BthreeVsAepN{} for AEP-full; per capability class the two
systems track each other: \BthreeAmbAuth{} against \AepAmbAuth{} on \textsc{auth},
\BthreeAmbPosOnly{} against \AepAmbPosOnly{} on \textsc{pos-only},
\BthreeAmbNoReadback{} against \AepAmbNoReadback{} on \textsc{none}.

For ambiguity, B3 minus AEP-full is \BthreeVsAepAmbDiffPP{} percentage points.
The stratified run-cluster \BthreeVsAepAmbDiffConfidence{}\% interval is
[\BthreeVsAepAmbDiffLow{},\,\BthreeVsAepAmbDiffHigh{}] percentage points,
which lies within a $\pm\BthreeVsAepAmbMargin{}$ percentage-point margin. We
introduce that margin in this revision; it was not preregistered, and we state
it as a stipulation rather than derive it. We do not claim it is immaterial to
an operator: \cref{sec:threats} records that we have not evaluated declared
ambiguity as an operational outcome, and that the artifact has no escalation
mechanism, so we have no basis on which to price one. The result
supports equivalence under that stated, post hoc margin, but only for this
fixed design. Each arm has \BthreeVsAepAmbClusters{} run clusters distributed
over \BthreeVsAepAmbStrata{} matched strata---three runs per stratum---so the
cluster-aware interval is a sensitivity analysis, not a strong general
equivalence result. The execution-level Fisher value
($p = \BthreeVsAepAmbP{}$) is not used as equivalence evidence.

\textbf{A word on what the zeros can and cannot support.} Fisher's exact test
on $\BthreeVsAepDupCount{}/\BthreeVsAepN{}$ against
$\BthreeVsAepDupCount{}/\BthreeVsAepN{}$ returns
$p = \BthreeVsAepDupP{}$ for duplicates and $p = \BthreeVsAepLostP{}$ for lost
effects. Neither $p$-value carries information: a test between two zero
counts has no power, and failing to distinguish them is not evidence that they
are the same. What two zeros do support is a \emph{bound}. The unit for that bound
is the run, not the execution: each arm's \BthreeVsAepN{} executions are
\BthreeVsAepRuns{} runs of ten, and ten executions that share a run, a crash
point and a system are not ten independent trials. The one-sided
95\% Wilson upper confidence bound on a zero numerator over
\BthreeVsAepRuns{} run clusters is \AblationZeroUpperRun{}\%. This is a genuinely
one-sided bound using the 95th normal quantile; it is not the upper endpoint of
a two-sided 95\% Wilson interval, which uses the 97.5th quantile. Each system's
rate is individually bounded this way. When the two system bounds are stated
simultaneously, Bonferroni gives joint coverage of at least 90\%, not 95\%.
Under that simultaneous event, both rates lie in
$[0,\AblationZeroUpperRun{}\%]$, so their absolute difference is at most
\AblationZeroUpperRun{} percentage points. This finite-sample bound is the support
for ``no observed difference'' in the two zero-event metrics; the corresponding
Fisher values are not evidence of equivalence.

Treating the executions as independent would give
\AblationZeroUpperExec{}\%. We use the run, for the same reason the Fisher
values in \cref{tab:ablation} are not read as equivalence: with no events
observed, nothing in these data estimates the within-run correlation that
execution-level independence would require. The choice of unit decides nothing
claimed below:
the baseline duplicate rates this bound is contrasted against are
\BaselineDupLowPct{}--\BaselineDupHighPct{}\%, more than an order of magnitude
from either figure, so the conservative unit is taken because it is the
defensible one and not because the comparison needs it.

\textbf{We report this as a finding rather than as a limitation, because it
reassigns our own headline claim.} The property that distinguishes AEP from
B0, B1 and B2 --- and the property the paper is named for --- is produced by
the pre-dispatch record and by the transition table that refuses to
re-enter \textsc{about-to-fire}. It is not produced by waiting for
\texttt{fsync}. A reader who wants the detection guarantee and not the
barrier's cost can have it, and \cref{sec:eval-deployment} prices that
configuration explicitly.

This also sharpens what B4 shows. B4 has a durable record, acknowledged by the
same barrier, and duplicates at the highest rate in the study; B3 has a durable
record, no barrier, and duplicates at zero. Durability is neither necessary nor
sufficient for detection. What is necessary is that the record be written
\emph{before} the call and that no path re-enter the dispatching state.

\subsubsection{Prevention is the barrier's, and it needs the right fault}
\label{sec:eval-prevention}

If detection does not need the barrier, what does the barrier do? It withholds
the external call when the durability of the intent can no longer be confirmed.
That is a different guarantee, it is worth stating in its own metric, and no
worker-crash fault can exercise it --- which is why \cref{sec:eval-detection}
returns a null. The fault that exercises it is loss of Redis itself.

The metric is the \emph{unwanted-applied-effect rate}: the fraction of
executions in which a real, non-idempotent mutation reached the provider while
the durability of its own intent record could no longer be established. It is
not the same quantity as an undetected duplicate and it is not measured on the
same regime; keeping the two apart is the whole point of this subsection.

We hard-kill Redis at \texttt{after\_intent\_before\_barrier}: the instruction
boundary between the intent CAS and the acknowledgement. AEP-full is inside
\texttt{WAITAOF} at that moment; B3 is not, because B3's ablation is precisely
that round trip. One execution per run, no worker crash --- the fault under
study is Redis dying.

\paragraph{The controlled fault} The result below is collected under a
\emph{controlled} injector, and that is the change that makes it worth
reporting. \texttt{docker kill} alone lands \UncontrolledLanding{}\,ms after it
is issued on this host; freezing the container first with \texttt{docker pause}
and killing it afterwards lands in \ControlledLanding{}\,ms, with a third of the
spread, and leaves the container in the same end state. The fault class does not
change --- it is the same F3 crash-stop at the same instruction boundary --- so
the two are comparable, and the injector's own contribution to the race is
narrowed by measurement rather than argued away.

Over \ArmASessions{} sessions and all \ArmAClasses{} capability classes,
AEP-full's unwanted-applied rate is \ArmAAepRate{} (\ArmAAepApplied{} of
\ArmAAepN{} executions) and B3's is \ArmABthreeRate{}
(\ArmABthreeApplied{}/\ArmABthreeN{}), on the same injected fault. The
between-session spread in AEP-full's count --- the quantity a single collection
cannot see, and the one the earlier version of this result was criticised for
not reporting --- collapses to \ArmASpreadMax{} count per class across the
\ArmASessions{} sessions. B3 stays at its ceiling throughout, never below
\ArmABthreeFloor{} of \AepKillRuns{} in any class in any session, which is what
rules out the reading that the injector simply disabled both arms.

\paragraph{The result fell below its own pre-registered bound, and we report
that as a surprise rather than a success} The prediction was registered before
the data existed: \ArmAPredicted{}, bounded [\ArmABoundLow{},
\ArmABoundHigh{}], derived from the measured landing distribution as the
fraction of a $1\,000$\,ms \texttt{WAITAOF} window the injector leaves open.
The observed \ArmAAepRate{} is \emph{below} that bound. The pre-registration
fixed in advance what a rate outside the bound would mean --- the protocol
behaving unexpectedly, not the mechanism succeeding --- and we honour that here
rather than presenting a lower-than-predicted number as a better result.

The reason the bound was wrong is measurable, and it is not calibration: the
model behind it fails on two independent premises that push in opposite
directions. First, it treats the exposure window as the injector's bench-measured
landing latency. That does not survive contact with a collection:
\ModelAcksPastLanding{} of \ModelAckN{} uncontrolled AEP-full runs received a
genuine \texttt{WAITAOF} acknowledgement \emph{after} the bench landing had
supposedly passed, and a \texttt{WAITAOF} acknowledgement requires a live
server. Timed during the collections themselves, the same injector call takes
\ModelInSituKill{}\,ms against the bench's \ModelBenchKill{}\,ms ---
\ModelInSituRatio{}$\times$, with ranges that do not overlap. Second, it treats
the \texttt{WAITAOF} wait as uniform on $U(0, 1\,000)$\,ms. Measured over the
runs the model was calibrated against, the fraction of waits under
\ControlledLanding{}\,ms is \ModelEmpiricalWindow{}, where uniformity assigns
\ModelUniformWindow{} --- \ModelUniformOverstates{}$\times$ too generous, and
worse the narrower the window. The two errors partly cancelled at the
uncontrolled window and could not at the controlled one.

The calibration that appeared to license the bound compared the wrong two
quantities. The model predicts \emph{dispatch}; the estimand is an
\emph{applied effect}, and dispatch is necessary for one without being
sufficient. On the same runs the dispatch rate is \ModelDispatchRate{} and the
applied rate is \ModelAppliedRate{}. Against the quantity it actually predicts,
the model was already wrong on its own calibration data. We record this because
a pre-registered bound that agrees for the wrong reason is worse than one that
disagrees, and only the second kind is visible without looking.

\paragraph{What this does and does not license us to say} We report a rate and
its mechanism, not a guarantee. \textbf{We do not claim that AEP-full prevents
the unwanted effect.} The residual is non-zero and is irreducible under this
fault class: freezing Redis synchronously at the checkpoint would make
\texttt{WAITAOF} unanswerable, but B3 reaches the same checkpoint and still
needs two further Redis calls, so the same freeze would stop B3 dispatching too
and the contrast would vanish because the injector had disabled both arms. What
the controlled fault buys is that the residual is now bounded by something we
measured rather than by this host's \texttt{docker} latency --- and, as above,
the bound we derived from that measurement was itself wrong in a way the data
can show. The magnitude is not a constant of the protocol and this paper does
not offer it as one.

\paragraph{The uncontrolled cell replicates the direction, and is not pooled
with the controlled one} The original collection used \texttt{docker kill}
alone, and is kept as the uncontrolled replication. It is reported separately
rather than pooled: the two injectors have different landing distributions, and
combining them would trade a stated limitation for an unstated confound. There
the rate is \AepUnwantedRate{} for AEP-full against \BthreeUnwantedRate{} for B3
(Fisher's exact test, two-tailed, $p = \UnwantedP{}$), a difference of
\UnwantedPrevented{} executions on the same injected fault --- the same
direction as the controlled cell, at a magnitude the uncontrolled injector sets.
This is the only fault in the study under which the two systems diverge;
everywhere else they differ by at most $\BthreeVsAepAmbDelta{}$ executions in
\BthreeVsAepN{} (\cref{tab:ablation}). Both systems declared ambiguity on all
\AepKillRuns{} runs: the endpoint is \textsc{no-readback}, so neither can
resolve what it does not know, and both fail closed. The difference between them
is not in what they \emph{say} --- that is the detection guarantee, and
\cref{sec:eval-detection} already showed it does not depend on the barrier ---
it is in what they did to the outside world before saying it.

\textbf{The direction replicates; the magnitude moves, which is why the
controlled cell exists.} The cell was re-collected \ReplicationSessions{}
further times under an identical configuration, and AEP-full's unwanted-applied
count ranged \ReplicationAepMin{}--\ReplicationAepMax{} out of \AepKillRuns{}
while B3 stayed at \ReplicationBthreeApplied{}/\AepKillRuns{} in every one.
Taking the session as the unit, the difference averaged
\ReplicationPreventedMean{} effects, 95\% interval
[\ReplicationPreventedLow{},\,\ReplicationPreventedHigh{}]; against that spread
the controlled cell's \ArmASpreadMax{}-count per-class spread over
\ArmASessions{} sessions is the change worth reporting. Two negative findings
travel with the cell and are not softened by being moved: attributing the
movement to kill latency did not succeed at this design's precision, and B3 is
not the negative control an earlier version of this paper reported it to be ---
its own session-level interval is far too wide to have contradicted the
mechanism whatever it had done. The per-session tables carrying both --- the
uncontrolled counts, the four session-level kill-latency differences with their
intervals, and the same analysis for B3 --- are in the supplementary material.
The scope of the cell is one pre-acknowledgement crash point and one host, and
the \textsc{authoritative-readback} class was subsequently collected over
\ClassSessions{} sessions (\cref{sec:threats}) at a precision inadequate to the
question.

\paragraph{Where the fault lands after the branch point, the arms tie} The
same fault delivered \emph{after} dispatch separates nothing, and we collected
that rather than inferring it. Arming the kill at the last instruction before
transmission and delivering it $200$\,ms later puts it in flight, after both
arms have already dispatched. Over \InflightSessions{} sessions the two arms
differ by at most \InflightMaxDiff{} count in any capability class, and both sit
at their ceiling. Two qualifications belong with that number. The second session
is a deterministic replay of the first rather than an independent draw --- the
harness assigns seeds per cell and repetition, so where no race remains the
second collection reproduces the first exactly, which is itself the clearest
evidence that no race remains. And B3's lowest cell is
\InflightBthreeFloor{}/\AepKillRuns{}, exactly the floor below which the result
would have been read as a mis-timed injector rather than a tie; it is the
boundary, not headroom.

What is structural, across all three faults, is that the race exists for
AEP-full at all and cannot exist for B3.

\subsubsection{The barrier's durability claim, under the fault that can
actually test it}
\label{sec:eval-durability}

There remains a third thing the barrier is supposed to do, and it is the one
most readers assume it does first: protect the write-ahead record from being
lost when Redis dies. We tested that premise before building on it,
and for the fault everyone reaches for first, it is false.

\ProcessKillTrials{} trials, each phase-aligning the \texttt{everysec} timer
with a \texttt{WAITAOF} so the window under test is as wide as it can be,
writing a key \emph{without} waiting for it, hard-killing the container with
\texttt{docker kill -s KILL}, restarting, and reading back. Every kill landed
(every container came back with its \texttt{uptime} at zero) and every kill
landed inside the window (write-to-death
\ProcessKillWindowMin{}--\ProcessKillWindowMax{}\,ms against a $1\,000$\,ms
fsync period).
\textbf{Not one unfsynced write was lost: $\ProcessKillUnackLost{}$.}
% source: reports/raw/e1-durability-window.txt; Session 3B section C.2.

The mechanism is not subtle once seen. \texttt{appendfsync everysec} defers the
\texttt{fsync(2)}, \textbf{not} the \texttt{write(2)}: Redis writes the AOF
buffer to the operating system on every event-loop iteration. A
\texttt{SIGKILL} destroys the \emph{process}; the kernel and its page cache are
untouched, and a kernel that is still running flushes those bytes to disk. The
same result reappears at $n{=}\KillCanaryN{}$ inside the ablation cells above, where a
canary write made immediately before each kill survived in \AepKillCanary{} and
\BthreeKillCanary{} runs.

That narrows the claim to a named fault class --- loss of the page cache: host
power failure, kernel panic, VM destruction --- and an earlier draft of this
paper stopped there, naming a fault it had not injected. We now inject it.

\paragraph{Making the storage lose the writes} The probe is the one above with
the fault swapped. A loop-backed block device carries a \texttt{dm-flakey}
target whose \texttt{drop\_writes} feature silently discards every write bio
while serving reads correctly, which is what a storage stack looks like to a
kernel whose writes will never reach the platter. Redis 7.2.5 --- the same
build as the rest of the evaluation --- runs with its append-only file on an
ext4 filesystem over that device. Each trial writes one key through
\texttt{WAITAOF} and one without, switches the device into
\texttt{drop\_writes}, kills the server, unmounts to discard the page cache,
restores the device and reads both keys back. The switch is the instant of
power loss: what was \texttt{fsync}-ed before it is on the device, what was
still dirty is not.

Two details decide whether this measures anything. \texttt{dmsetup suspend}
freezes the filesystem and flushes outstanding I/O by default, either of which
would push the unacknowledged write out before the fault took effect and
produce a confident null; the reload runs with \texttt{--nolockfs --noflush}.
And a self-test runs first on the device alone, with no Redis in the picture: a
file \texttt{fsync}-ed before the switch must survive and one \texttt{fsync}-ed
after it must not. The probe refuses to report Redis trials unless that holds.
% source: experiments/flakey_write_loss.py; reports/raw/g2-flakey-write-loss.txt.

\paragraph{The result} The probe ran \FlakeyN{} trials in \FlakeyReplications{}
independent replications. Each replication builds the device stack from
scratch, and each \emph{trial} within it rebuilds the filesystem with
\texttt{mkfs.ext4} and starts a fresh server, so no trial can read a key a
previous trial's append-only file still holds. \textbf{Separation was perfect in
every replication}: the acknowledged record survived
\FlakeyPerRepAckSurvived{} and the unacknowledged one was destroyed
\FlakeyPerRepUnackLost{} in each, \FlakeyAckSurvived{} and \FlakeyUnackLost{}
overall. Against the same two keys under a process kill,
$\ProcessKillUnackLost{}$ were lost; the fault classes separate completely.

We give no $p$ for either comparison, and the reason is that no unit supports
one: at the trial level the test would assume an independence the shared device
stack violates and would ignore the pairing --- both records are written in the
same trial and meet the same write-loss event --- while at the replication
level, three against three, the two-sided minimum a perfectly separated result
can reach is $0.1$, which is the design's floor rather than a statement about
these data. The complete separation in every replication is the finding, and it
does not need one. The filesystem was mounted
\texttt{ext4 rw,relatime} with its journal enabled and no \texttt{nobarrier},
and the server's \texttt{appendfsync} was read back as \texttt{everysec} from
the running instance; the probe refuses to measure if either is otherwise,
because under \texttt{always} there is no unacknowledged write to lose and
under \texttt{nobarrier} a surviving acknowledged one would prove nothing.

\textbf{The barrier's durability benefit to the \emph{record} is therefore real
and now measured, and the fault class matters more than the mechanism.} What
that benefit is worth to a \emph{dispatch} is the next subsubsection's question,
and the answer there is not the one we expected. The Redis documentation's
``you may lose 1 second of data if there is a disaster''~\cite{redis-persistence}
is correct; what our two probes add is that ``disaster'' must mean the host,
that the widely repeated reading of the sentence as covering a crashed Redis
\emph{process} is wrong, and that once the host is the fault the sentence is
exactly right.

One qualification, and it is the same one the process-kill probe carries: the
exposure window is deliberately maximised. The unacknowledged write is issued
\FlakeyWindowMin{}--\FlakeyWindowMax{}\,ms before the device stops accepting
writes, inside a $1\,000$\,ms fsync period that a preceding \texttt{WAITAOF}
has just reset. \FlakeyUnackLost{} is therefore the loss probability for a
power cut placed where it does the most damage, not the rate a uniformly-timed
one would produce; that rate is bounded above by it and approaches it as the
write approaches the start of the period. The design is deliberate and it is
the same design that made the process-kill probe's $\ProcessKillUnackLost{}$
strong: both probes
give the loss hypothesis its best case, and only one of them finds loss.

\subsubsection{Running the protocol under that fault, and what it refuted}
\label{sec:eval-writeloss-cell}

The probe above tests \texttt{WAITAOF}, not AEP. It runs two Redis keys against
a device that stops accepting writes; it does not run AEP-full and B3 through
the harness under that fault, and so it measures no difference in \emph{protocol}
outcomes. What it establishes is the premise the durability argument rests on.
We then ran the extension, and it is reported here rather than as a caveat
because it refuted what we predicted.

\paragraph{What was registered, and what happened} We pre-registered that
AEP-full would withhold dispatch when its write-ahead record was destroyed and
that B3 would proceed. A full session was collected with the harness's Redis
running its append-only file on the flakey device: both arms,
\WriteLossRunsPerArm{} runs each, \WriteLossExecPerRun{} executions per run,
\texttt{drop\_writes} armed at the intent checkpoint, without a delivery
failure. \textbf{The prediction is refuted.} AEP-full applied an effect in
\WriteLossAepApplied{} of \WriteLossExecPerArm{} executions and B3 in
\WriteLossBthreeApplied{}; both arms finished at ceiling, and neither withheld.

\paragraph{Why, and it is the more useful finding} \textbf{\texttt{WAITAOF}
returned success for every durability acknowledgement requested after the device
had stopped accepting writes} --- \WriteLossAcksAfterFault{} of them, with
\WriteLossAckFailures{} failures. \texttt{dm-flakey}'s \texttt{drop\_writes}
discards each write bio and reports success, so the kernel sees no error,
\texttt{fsync(2)} succeeds, and Redis has no way to learn otherwise. The barrier
withholds dispatch on a \emph{failed or absent} acknowledgement; it was handed a
successful one and dispatched, behaving exactly as specified on an input that
was false.

\paragraph{The bounds, stated narrowly} \textbf{This is not evidence that the
barrier fails to withhold under write loss}: the condition the experiment exists
to create --- the barrier meeting an acknowledgement that fails or never arrives
--- was never reached, so no run exercised withholding. \textbf{Neither is it
evidence that the barrier works.} What it establishes is that the durability
guarantee rests on an unstated premise --- that a \texttt{WAITAOF} success means
the record reached the device --- and that under storage which discards writes
\emph{silently} the premise is false and the protocol cannot detect it. Storage
returning \texttt{EIO} would let the acknowledgement fail and put the barrier
back on the input its guarantee is defined over, so this result bounds the
guarantee to storage that reports its own failures and says nothing about
storage that does not. It is evidence about storage, not about the protocol. One
host, one Redis version (7.2.5, \texttt{appendfsync everysec}), and we claim it
for no other.

\subsection{RQ3 --- cost, and how much of it is optional}
\label{sec:eval-rq3}

Only crash-free runs on a suspend-disabled host contribute. The provider is
configured with a constant $2\,000$\,ms delay, so one round trip is the floor
and the interesting quantity is the increment over it.

% GENERATED by scripts/paper_tables.py -- do not edit.
% Source: analysis/latency-and-throughput.csv
% 'median' is step_latency_ms_median over crash-free runs only
% (denominator = overhead_runs_crash_free, printed).
% 'over B0' subtracts B0's median (2010.165771 ms), which is
% the no-protocol system on the same host and the same provider
% (configured delay 2000 ms; B0 sits above it by
% 10.2 ms, which is harness cost, not protocol cost).
% E5 gate: a run contributes a duration only if the host was declared
% non-suspending BEFORE the run and no wall-versus-monotonic
% divergence was observed in it. Counts are unaffected.
\begin{table}[t]
\centering
\caption{Median step latency, crash-free runs only, E5-gated. Increments are over B0, the no-protocol system on the same host and provider. The whole write-ahead protocol minus the barrier is the B3 row; the barrier is the difference between B3 and AEP-full.}
\label{tab:latency}
\small
\begin{tabular}{@{}lrrr@{}}
\toprule
system & runs & median step (ms) & over B0 (ms)\\
\midrule
B0 no protocol & 3 & 2\,010.2 & ---\\
B1 $+$ lease & 3 & 2\,014.0 & 3.9\\
B2 $+$ fenced CAS & 3 & 2\,017.4 & 7.3\\
B3 full protocol, no barrier & 3 & 2\,038.2 & 28.0\\
\textbf{AEP-full} & 3 & 4\,004.9 & 1\,994.7\\
B4b durable, 1 attempt & 3 & 4\,013.4 & 2\,003.3\\
B4 durable, $\infty$ attempts & 3 & 4\,015.7 & 2\,005.5\\
\bottomrule
\end{tabular}
\end{table}


\textbf{The overhead is the barrier, and the rest is not separable from
zero.} A median step costs \AepStepMedian{}\,ms under AEP-full against
\BzeroStepMedian{}\,ms under B0, which does nothing but call the provider.
Almost all of that gap is one mechanism. At fifteen crash-free runs per arm the
two \texttt{WAITAOF} round trips cost \BarrierCostFifteen{}\,ms between them,
with a 95\% interval of
[\BarrierCostFifteenLow{}, \BarrierCostFifteenHigh{}]\,ms.
% generated/numbers.tex | ws5-2026-09-10/t1-p0-everysec/analysis/
% per-execution.csv, cluster bootstrap over runs, 10000 resamples, seed
% 20260806: median AEP-full - median B3.

Everything the write-ahead protocol does apart from waiting for \texttt{fsync}
--- the intent ledger, the vault, the request binding, the preflight, the
recovery service --- costs
\ProtocolMinusBarrierFifteen{}\,ms
[\ProtocolMinusBarrierFifteenLow{}, \ProtocolMinusBarrierFifteenHigh{}]\,ms on
a two-second call. That interval spans zero. B3's arm is a mixture of two
modes, and restricting both arms to their lower mode gives
\ProtocolMinusBarrierLowerMode{}\,ms
[\ProtocolMinusBarrierLowerModeLow{}, \ProtocolMinusBarrierLowerModeHigh{}]\,ms,
which spans zero as well. We report both because they were both pre-registered
and they disagree, and choosing between them after seeing them is the move this
paper is arguing against.
% generated/numbers.tex | same CSV and bootstrap: median B3 - median B0,
% pooled over both modes, and again restricted to each arm's lower mode. The
% lower-mode boundary is the splitter's own lower_max, fixed before this arm
% was collected.

\textbf{So the decomposition supports an ordering and not a partition.} The
barrier dominates: its interval is nowhere near zero and the residual's
contains it under both readings. What cannot be said is what fraction of the
protocol's cost is optional, because at this precision the non-barrier
remainder is not distinguishable from nothing at all. An earlier draft of this
paper quoted that remainder as a point figure of about thirty milliseconds; it
rested on three runs per arm, and twelve more did not confirm it.

Read against \cref{sec:eval-detection} the ordering still says something
stronger than ``AEP costs two seconds''. Whatever the detection guarantee
costs, it is the difference between B0 and B3 --- B3 is the system that already
has that guarantee in full --- and it is small enough that fifteen runs cannot
separate it from zero. The \BarrierCostFifteen{}\,ms buys the prevention
guarantee of \cref{sec:eval-prevention} and nothing else.

That figure is also a property of a \emph{configuration}, not of AEP. Under
\texttt{appendfsync everysec} an acknowledgement cannot return until the next
scheduled fsync, and those are one second apart, so a barrier costs up to a
second and the protocol pays that twice. The measured total sits near the top
of that range rather than in the middle of it; the reason is unmeasured.

\subsubsection{The barrier is a deployment choice, with three measured points}
\label{sec:eval-deployment}

The $\approx 2$-second figure above is the one a reader would carry away as
``AEP's overhead''. It is not: it is the price of one of the protocol's two
guarantees, under one of the durability settings that guarantee can be bought
at. \Cref{tab:deployment} is the decision as a deployment actually faces it.

% GENERATED by scripts/paper_tables.py -- do not edit.
% Sources: analysis/latency-and-throughput.csv (everysec rows),
%          fsync-always/analysis/latency-and-throughput.csv (always).
% step_latency_ms_median over crash-free, E5-gated runs only.
% Barrier column = row median - the B3 median collected under the
% SAME fsync policy (everysec 2\,038.2 ms, always 2\,048.4 ms),
% i.e. the same protocol with the barrier ablated and nothing else
% changed. Never across policies: that would assume the ablated
% protocol's own writes cost the same under both, which is exactly
% what the two B3 numbers are here to establish rather than assume.
\begin{table*}[t]
\centering
\caption{The barrier is a deployment choice, not a fixed cost. All three rows run the same pre-dispatch intent ledger and therefore make the same detection claim; they differ in what they pay to also withhold dispatch when durability cannot be confirmed. The barrier column is each row's own median minus a B3 median collected under the same \texttt{appendfsync} policy. Crash-free runs only, E5-gated, 2\,000\,ms provider floor.}
\label{tab:deployment}
\small
\renewcommand{\arraystretch}{1.08}
\begin{tabularx}{\textwidth}{@{}>{\raggedright\arraybackslash}Xrrrc>{\raggedright\arraybackslash}p{0.20\textwidth}@{}}
\toprule
configuration & \shortstack{median\\(ms)} & \shortstack{over\\floor} & \shortstack{barrier\\(ms)} & prevents & claim\\
\midrule
AEP-full, \texttt{everysec} & 4\,004.9 & 2\,004.9 & 1\,966.7 & yes & detection + prevention\\
AEP-full, \texttt{always} & 2\,063.4 & 63.4 & 15.0 & yes & detection + prevention\\
B3-mode, \texttt{everysec} & 2\,038.2 & 38.2 & 0.0 & no & detection only\\
\midrule
\multicolumn{6}{@{}p{0.96\textwidth}@{}}{\footnotesize The detection claim of \cref{tab:outcomes} shows no observed difference down the whole table: it is produced by the pre-dispatch record plus no re-entry, which all three rows have, and \cref{tab:ablation} is the ablation that shows it. What the barrier buys is the last column's second word, and \cref{tab:killablation} is what it is worth. `Over floor' is the same median less the provider's 2\,000\,ms delay, and so includes the 28.0\,ms the protocol costs with the barrier already removed.}\\
\bottomrule
\end{tabularx}
\end{table*}


Three points, and the third is the one the earlier framing of this paper
concealed.

\paragraph{\texttt{everysec} with the barrier} The default, and the expensive
one. An acknowledgement must wait for the next scheduled fsync, up to a second
away, and the protocol waits twice.

\paragraph{\texttt{always} with the barrier} The identical crash-free cells ---
same systems, endpoint, regime, seeds, host and provider --- against a Redis
whose only changed line is \texttt{appendfsync always}.
% source: reports/raw/f0iii-appendfsync-always.txt and the G2 re-run that
% added the B3 arm; both refuse to proceed unless CONFIG GET appendfsync
% reads back 'always' on the live server.
The barrier's cost is \BarrierCostAlwaysFortyFive{}\,ms
[\BarrierCostAlwaysFortyFiveLow{}, \BarrierCostAlwaysFortyFiveHigh{}]\,ms over
fifteen runs per arm: negative, and not distinguishable from zero. Under
\texttt{always} the fsync has already happened by the time the write returns,
so \texttt{WAITAOF} has nothing left to wait for.
% generated/numbers.tex | fsync-always-2026-09-14/analysis/
% latency-and-throughput.csv and per-execution.csv, same bootstrap as above.

\textbf{That is the measurement this section turns on.} The barrier does not
become cheap under \texttt{always}; it costs nothing, and the
\BarrierCostFifteen{}\,ms it costs under \texttt{everysec} is therefore the
price of the fsync boundary rather than of the barrier. A protocol step under
\texttt{everysec} waits for the next scheduled flush, and there are two of
them. Nothing in the mechanism is expensive.

The sign is not a puzzle to explain away. The two medians it separates are
\AepAlwaysFortyFiveMedian{}\,ms with the barrier and
\BthreeAlwaysFortyFiveMedian{}\,ms without it, on a two-second call; the gap
between them is noise around a mechanism that is doing no waiting. This arm was
pre-registered with the interval's exclusion of zero explicitly ruled out as an
acceptance criterion, because a mechanism with nothing to wait for should cost
nothing and a measurement of nothing lands on either side of it.
% generated/numbers.tex | fsync-always-2026-09-14/analysis/
% latency-and-throughput.csv, step_latency_ms_median per system.

Both figures are ablation differences \emph{within} one policy, which required
collecting the barrier-less arm twice. It would have been easier to subtract
the \texttt{everysec} B3 median from the \texttt{always} AEP-full median, and
that would have assumed the ablated protocol's own writes cost the same under
both policies. They nearly do --- B3 is \BthreeStepMedian{}\,ms under
\texttt{everysec} and \BthreeAlwaysFortyFiveMedian{}\,ms under
\texttt{always} --- but
``nearly'' is a measurement, and an order-of-magnitude claim resting on an
unmeasured premise is the kind of number this paper is trying not to publish.

\paragraph{How precise these two figures are} Each is a difference of two
medians over fifteen runs per arm, bootstrapped over runs rather than
executions. Under \texttt{everysec} the barrier costs
\BarrierCostFifteen{}\,ms, [\BarrierCostFifteenLow{}, \BarrierCostFifteenHigh{}];
under \texttt{always} it is \BarrierCostAlwaysFortyFive{}\,ms,
[\BarrierCostAlwaysFortyFiveLow{}, \BarrierCostAlwaysFortyFiveHigh{}]. The first
excludes zero by a wide margin and the second contains it, which is the whole
of the claim: the policy decides whether the barrier is expensive, and the
barrier itself is not.

\paragraph{B3-mode: the barrier removed entirely} Not a baseline we are beating
but a supported configuration of the same implementation, licensed by
\cref{sec:eval-detection}. What it pays over a system with no protocol at all
is not resolvable from zero at fifteen runs per arm
(\ProtocolMinusBarrierFifteen{}\,ms,
[\ProtocolMinusBarrierFifteenLow{}, \ProtocolMinusBarrierFifteenHigh{}]). On the detection metrics, pooled across capability
classes, we observe no difference between it and AEP-full, supported by a bound
rather than a test: both the undetected-duplicate and lost-effect rates lie in
$[0,\AblationZeroUpperRun{}\%]$ at joint coverage of at least 90\%, and the
declared-ambiguity difference lies inside the $\pm\BthreeVsAepAmbMargin{}$
percentage-point margin \cref{sec:eval-detection} states as a post hoc
stipulation. \emph{We do not make that claim per capability class, and the
reason is a number rather than an oversight}: six simultaneous bounds instead of
two, over a third as many run clusters, puts the Bonferroni width at
\AblationZeroUpperPerClass{} percentage points, and a band that wide is not a
detection guarantee. What B3-mode forfeits is \cref{sec:eval-prevention}: under
a hard Redis kill it will put an effect on the wire that AEP-full withholds, at
\BthreeUnwantedRate{} against \AepUnwantedRate{}, and its intent record will not
survive a host-level write loss (\cref{sec:eval-durability}). That second
clause is about the \emph{record}, and it does not license the protocol-level
inference a reader may draw from it: when we ran both arms under exactly that
fault, neither withheld, because the storage discarded the writes while
reporting success and the barrier was handed an acknowledgement that looked
good (\cref{sec:eval-writeloss-cell}).

\textbf{So the question a deployment answers is not ``can we afford AEP'' but
``do we need prevention, and at which fsync policy''.} If losing Redis is the
fault to worry about --- an instance on the same host as the workers, an
unreplicated instance, a VM that can vanish --- the barrier is the point of the
protocol and \texttt{always} makes it nearly free at this workload's shape. If
Redis is well protected and worker crashes are the concern, B3-mode gives the
paper's headline result for a cost this evaluation cannot distinguish from zero.

\textbf{One measurement this table does not support} is what \texttt{always}
costs under load. The Redis
documentation calls it ``very very slow, very safe''~\cite{redis-persistence}, a
judgement about write throughput that our workload cannot exercise: two workers,
a two-second provider delay and a handful of writes per step leave the system
latency-bound and nowhere near an fsync-rate limit. Our throughput figure is in
fact \emph{higher} under \texttt{always} ($\AepThroughputAlways{}$ versus
$\AepThroughputEverysec{}$ executions/s), which is a statement about this
workload's shape and not about Redis.

\subsection{RQ4 --- recovery}
\label{sec:eval-rq4}

Recovery behaviour is reported through the outcome distribution rather than a
success rate, for a reason worth stating: ``recovery success'' means two
different things either side of whether a recovery service was running. Where
one ran, a crashed execution reaching a terminal classification was
\emph{recovered}. Where none ran --- B0, B1, B2, B4, B4b declare none --- the
same execution reached its classification because a supervisor ran the step
again, which is not recovery but the thing recovery exists to avoid. Our
analysis excludes runs without a recovery service from that metric rather than
scoring them zero or one, because both would be statements about a component
that was not present.

Within that scope: every AEP-full and B3 execution in the crashed regime
reached a terminal classification, and across \RunsCollected{} runs exactly
one recorded a reconciliation disagreement --- every other completed run's
event log agreed with its independently written oracle ledger. The harness
performs this reconciliation after each completed run; it checks definitive
applied/non-applied assertions, effect-count bounds, attribution, and duplicate
predictions rather than trusting a line count.
% source: analysis/coverage.json and the per-run reconcile step.

That one is worth stating precisely, because the check exists to be believed.
It was a B4b cell in which the system classified all ten executions
\textsc{confirmed-applied} while the provider's ledger recorded two, eight
disagreements of the form ``the system asserted the effect was applied and the
ledger records none''. Both sibling runs of the same cell recorded ten, the
provider logged no error, and its run-log held three lines against eleven for
each sibling --- so the provider recorded almost nothing while answering the
workers successfully. We could not determine the cause. A re-run of the
identical cell, same seed and same configuration, agreed.

We treat it as a \emph{void} run on the rule the durability probe uses: a
trial whose measurement precondition failed says nothing about the system in
either direction. Counting it would let a broken ground truth set a baseline's
rate; discarding it quietly and keeping the re-run would be re-running until
the number looks right. So the re-run is what \cref{tab:outcomes} contains,
the voided attempt must ship in the external raw archive under
\texttt{results/voided/} with this explanation beside it; it is not present in
the current Git repository. The count above is one rather than zero.
Recovery latency is not reported as an absolute figure: the E5 gate leaves too
few gated crashed runs to support a distribution, and \cref{sec:threats}
records this as a gap rather than filling it with ungated numbers.
